
\section{第 4、5 章概述}
\label{sec:4-1}

第 3 章已经为相对 spin 拉格朗日子流形 $L$ 构造了相伴的滤过
$A_\infty$ 代数 $(C(L;\Lambda_{0,\mathrm{nov}}),\mathfrak m)$，也为相对
spin 对 $(L^{(0)},L^{(1)})$ 构造了滤过 $A_\infty$ 双模
$C(L^{(1)},L^{(0)};\Lambda_{0,\mathrm{nov}})$。本章和下一章考察这个滤过
$A_\infty$ 代数、这个滤过 $A_\infty$ 双模以及弗洛尔上同调对哈密顿形变和
定义中其他参数的依赖性；这些参数包括相容殆复结构 $J$、扰动（多值截面）
$\mathfrak s$、$\mathfrak s^{-1}(0)$ 的三角剖分，以及光滑奇异单形可数集的
选取。我们还将在定界余链 $b$（即 Maurer--Cartan 方程的解）所成的集合上
定义规范等价关系，并研究弗洛尔上同调对定界余链 $b$ 的选取有何依赖。在这
两章中，始终假定拉格朗日子流形 $L$ 可定向且相对 spin，并假定任意所论的
拉格朗日子流形对都是相对 spin 的（见第 3.1 节的定义 3.1.1）。

首先说明弗洛尔上同调为何会依赖各种选择，尤其会依赖定界余链 $b$。回忆，
本书定义拉格朗日子流形的弗洛尔上同调时所面对的根本困难，是伪全纯圆盘模
空间一般具有余维一边界。这意味着，伪全纯圆盘模空间的基本闭链未必能按通常
意义良好定义。

唐纳森在研究满足 $b_2^+=1$ 的四维流形上反自对偶（ASD）方程解的模空间时，
发现了类似现象 \textbf{[Don86II]}。他注意到，若考察一个单参数扰动族，尤其是
黎曼度量的扰动族，模空间的同调类可能会在某个参数值处发生跳跃；这种现象称
为\emph{穿墙}。为简单起见，回忆模空间虚维数为零时穿墙如何发生。模空间带符号
计数所得的数就是相应不变量（唐纳森不变量）。给定一个单参数扰动族，就得到
一族模空间 $\mathcal M_s$（$s\in[0,1]$），其并
\[
  \bigcup_s\{s\}\times\mathcal M_s
\]
是一个一维流形。若对一般的单参数族，这个参数化模空间是紧的，那么它的边界
就是 $-\mathcal M_0\cup\mathcal M_1$，从而可以用配边论证证明该不变量与扰动
无关。然而，当冒泡现象的余维为一时，即使对一般的单参数族也无法完全避开
冒泡，所以这个参数化模空间一般并不紧。于是，配边论证不能证明这样构造的不
变量与所用扰动无关。

在规范场论中，唐纳森引入“腔室结构”来处理这一困难。在那里，困难有一个便于
控制的特点：可以用本质上只是上同调理论的阿贝尔规范场论来判定穿墙发生之处。
具体说，穿墙发生在存在非零反自对偶调和 $2$-形式、且其上同调类为整类的点。
特别地，若所取扰动不改变阿贝尔规范场论，尤其不改变黎曼度量，便不会发生
穿墙。因此，例如固定黎曼度量以后，唐纳森不变量就是良定义的。对满足
$b_2^+=1$ 的四维流形，塞伯格--威腾不变量也有类似问题。

在我们的情形中，问题更为复杂，因为穿墙可能出现的位置完全不受控制。因此，
为得到对扰动不变的对象，必须考察整个扰动族。为此，我们将研究（无障碍）
拉格朗日子流形模空间的无穷小结构。固定 $M$ 中一个拉格朗日子流形 $L$，并
固定一个相容殆复结构 $J$。我们要在拉格朗日子流形的扩展模空间中找出 $L$ 的
一个形式邻域。若 $L$ 附近另有拉格朗日子流形 $L'$，则规定 $L\sim L'$ 当且仅当
存在哈密顿同痕 $\psi_s$ 使 $L'=\psi_1(L)$。可以看出，这些拉格朗日子流形的
$\sim$-等价类集合，同 $H^1(L;\mathbb R)$ 中零点的一个邻域同构。（拉格朗日
子流形模空间的豪斯多夫性会带来一个问题；这里只关心局部问题，故不作讨论，
参见 \textbf{[Fuk02I]} 第 1 节。）

出于多方面的原因，尤其因为 $H_2(M;L)$ 中存在马斯洛夫指标非零的相对同调类，
我们希望把这个参数空间扩展到 $H^{\mathrm{odd}}(L;\mathbb R)$。在扩展后的空间
$H^{\mathrm{odd}}(L;\mathbb R)$ 中，只有 $H^1(L;\mathbb R)$ 方向具有直接的几何
意义；但仍可在其他方向上定义“弗洛尔（上）同调”的形变。正如下文将说明的，
我们的构造类似于用亏格零格罗莫夫--威腾不变量形变杯积的构造。设 $M$ 是辛流
形，空间 $H^2(M;\mathbb R)$ 局部参数化 $M$ 的辛结构模空间。格罗莫夫--威腾
理论以及由它得到的弗罗贝尼乌斯结构，可以看作由
$H^{\mathrm{even}}(M;\mathbb R)$ 参数化的 $M$ 上杯积形变族；这里
$H^{\mathrm{even}}(M;\mathbb R)$ 就是“$M$ 上辛结构的扩展模空间”。沿
$H^2(M;\mathbb R)$ 方向的形变具有明显的几何意义，即辛形式的形变；利用
格罗莫夫--威腾势，还能把形变空间扩展到所有 $H^{\mathrm{even}}(M;\mathbb R)$
方向。若把参数空间视为超空间，则还可以扩展到
$H^{\mathrm{odd}}(M;\mathbb R)$ 方向（见
\textbf{[ASKZ97, Dub96, KoMa94, Man99, Ran96]}）。

上面提到的拉格朗日子流形扩展模空间，是弗罗贝尼乌斯结构和格罗莫夫--威腾势
中出现的扩展模空间之自然类比。我们考察拉格朗日子流形 $L$ 在这个扩展模空间
中的形式邻域；它是模空间定义方程的形式幂级数展开之零点集，也就是仓西映射的
零点集。（这个形式幂级数是否收敛，目前尚不清楚。）换言之，我们把所研究的
模空间看成一个形式概形。上一章引入的定界余链集合
$\widehat{\mathcal M}(L)$，对本章将定义的适当规范等价取商后，恰好得到这个形式
概形 $\mathcal M(L)$；参见第 1.4、5.4 节。

格罗莫夫--威腾势的一个重要性质，是辛结构形变的扩展模空间不存在障碍；也就是
说，$H^{\mathrm{even}}(M;\mathbb R)$ 中零点的整个邻域都局部参数化这个扩展模
空间。原则上，$H^{\mathrm{odd}}(M;\mathbb R)$ 是障碍群；然而，对格罗莫夫--
威腾势而言，障碍映射或仓西映射恒等于零。这同镜像对称相容，具体如下。设
$M^\vee$ 是 $M$ 的镜像，则 $H^{\mathrm{odd}}(M;\mathbb C)$ 的一部分将同
$H^2(M^\vee;TM^\vee)$ 同构；后者正是 Kodaira--Spencer 形变理论的障碍所在的
向量空间。换言之，这是复结构形变理论的仓西映射所取值的空间。然而，对
卡拉比--丘流形，仓西映射自动为零。对经典方向
$H^1(M^\vee;TM^\vee)$，这一事实由 Bogomolov--田刚--Todorov 证明
\textbf{[Bog78, Tia87, Tod89]}；对其余扩展方向
$H^k(M^\vee;\Omega^{\mathrm{odd}-k}M^\vee)$，则由
Baranikov--孔采维奇证明 \textbf{[BaKo98]}。

另一方面，我们的拉格朗日子流形扩展模空间一般确实有障碍。相应的障碍映射或
仓西映射恰好就是我们的障碍类。因此，寻找（无障碍）拉格朗日子流形扩展模
空间的形式邻域，等价于研究障碍类在 $L$ 的一个邻域中的行为。要理解障碍类和
弗洛尔上同调对扰动的依赖，也就同这一研究密切相关。这里的障碍现象与前面提到
的穿墙现象直接相连：穿墙之所以发生，正是因为参数化的障碍类不为零。

下面说明第 4、5 章的内容。第 4 章研究滤过 $A_\infty$ 代数的同伦等价，第 5 章
研究滤过 $A_\infty$ 双模的同伦等价；两章都包含代数和几何两部分。代数部分是
$A_\infty$ 结构的同伦（同调）代数，已有多种参考文献（例如见
\textbf{[Fuk02II, 03II, Kel01, Lef03, MSS02, Smi00, Sei06]}）。不过，本书仍给出
自足的讨论，因为文献中尚未讨论滤过 $A_\infty$ 代数的情形；而且文献并没有
清楚说明它同定界余链规范等价之间的关系。为应用于 \textbf{[FOOO09I]}，我们还
需要讨论（同伦）幺元，并处理挠系数的情形。

第 4.2--4.5 节完全属于代数部分。（第 4.3.3 小节讨论
$\widehat{\mathcal M}_{\mathrm{weak,def}}(L)$ 上的规范等价。第 3 章引入 bulk 形变
时使用了由几何定义的算子 $\mathfrak q$；不过，只要任意滤过 $A_\infty$ 代数带
有一个满足定理 3.8.32 结论的算子 $\mathfrak q$，相应结果同样成立。）第 4.2 节
定义两个（滤过）$A_\infty$ 同态彼此同伦，并建立同态间同伦的基本性质；由此
可以定义两个（滤过）$A_\infty$ 代数同伦等价。定义 3.2.33 已经给出了（滤过）
$A_\infty$ 同态为弱同伦等价的概念。事实上，一个（滤过）$A_\infty$ 同态是
同伦等价，当且仅当它是弱同伦等价；这就是第 4.5 节所证明的定理 4.2.45。这个
结果给出了判定 $A_\infty$ 同态是否为同伦等价的便利准则。证明 $A_\infty$ 同态
之同伦具有基本性质所需的关键结果定理 4.2.34，则在第 4.4 节证明。

第 4.3 节定义两个定界余链规范等价，并建立其基本性质。这是一个等价关系；用它
对定界余链集合 $\widehat{\mathcal M}(C)$ 取商，得到集合
$\mathcal M(C)$。（第 5.4 节还会在 $\mathcal M(C)$ 上定义一个结构。）我们将证明，
$\mathcal M(C)$ 是滤过 $A_\infty$ 代数 $C$ 的同伦型不变量。

\liftIndexRangeStart{I140}
第 4.6 节转向几何情形，并证明与相对 spin 拉格朗日子流形 $L$ 相伴的滤过
$A_\infty$ 代数 $(C(L;\Lambda_{0,\mathrm{nov}}),\mathfrak m)$ 之同伦型（在第 4.3 节的意义下），
与构造中涉及的各种选择无关。更精确地说，我们将证明下面的定理 4.1.1。

设 $(M,\omega)$、$(M',\omega')$ 是辛流形，$L\subset M$、$L'\subset M'$ 是
拉格朗日子流形；又设 $\psi:M\to M'$ 是满足 $\psi(L)=L'$ 的辛微分同胚。分别
为 $M,L$ 和 $M',L'$ 固定构造
$(C(L;\Lambda_{0,\mathrm{nov}}),\mathfrak m)$ 与
$(C(L';\Lambda_{0,\mathrm{nov}}),\mathfrak m)$ 时所需的各种选择，即相容殆复
结构、扰动、三角剖分和光滑奇异单形可数集。

\begin{theorem}{4.1.1}\label{thm:4-1-1}
在上述情形中，$\psi$ 诱导一个同伦含幺滤过 $A_\infty$ 同态
\[
  \psi_*:
  (C(L;\Lambda_{0,\mathrm{nov}}),\mathfrak m)
  \longrightarrow
  (C(L';\Lambda_{0,\mathrm{nov}}),\mathfrak m),
\]
它是滤过 $A_\infty$ 代数的同伦等价。
\end{theorem}

取 $M=M'$、$L=L'$，定理 4.1.1 就说明滤过 $A_\infty$ 代数
$(C(L;\Lambda_{0,\mathrm{nov}}),\mathfrak m)$ 在同伦等价意义下与构造中所作的
各种选择无关。我们还将证明，定理 4.1.1 中的同伦等价 $\psi_*$ 在同伦意义下
也与各种选择无关；即有下述结论。

\begin{theorem}{4.1.2}\label{thm:4-1-2}
在定理 4.1.1 的情形中，设 $\{\psi_s\}_s$（$0\le s\le1$）是一族辛微分同胚，
且每个 $s$ 都满足 $\psi_s(L)=L'$。那么 $\psi_{0*}$ 与 $\psi_{1*}$ 作为同伦
含幺滤过 $A_\infty$ 同态彼此同伦。
\end{theorem}

这两个定理连同第 4.2--4.5 节的代数结果蕴含：定界余链规范等价类的集合
$\mathcal M(L)$ 在典范同构意义下与各种选择无关；它们还蕴含弗洛尔上同调
$HF((L,b_1),(L,b_0);\Lambda_{0,\mathrm{nov}})$ 只依赖于 $b_i$ 的规范等价类。
更精确地说，有下述结论。

\begin{theorem}{4.1.3}\label{thm:4-1-3}
在定理 4.1.1 的情形中，存在双射
\[
  \psi_*:\mathcal M(L)\longrightarrow\mathcal M(L'),
\]
它只依赖于 $\psi$ 的同伦类（这里的同伦取定理 4.1.2 的意义）。

进一步，若
$b_0,b_1\in\widehat{\mathcal M}(L)$、
$b'_0,b'_1\in\widehat{\mathcal M}(L')$，且
$\psi_*[b_i]=[b'_i]\in\mathcal M(L')$，则存在典范同构
\[
  \psi_*:
  HF((L,b_1),(L,b_0);\Lambda_{0,\mathrm{nov}})
  \longrightarrow
  HF((L',b'_1),(L',b'_0);\Lambda_{0,\mathrm{nov}}).
\]
\end{theorem}

引言中定理 B 的 (B.3)、(B.4) 可由定理 4.1.3 及其
`def,weak' 版本推出；见第 4.6.5 小节。

第 5 章处理有两个拉格朗日子流形或两个（滤过）$A_\infty$ 双模的情形。第 5.1
节讨论弗洛尔上同调理论中出现的各种诺维科夫环之间的关系；第 5.2 节仿照（滤
过）$A_\infty$ 代数的情形，研究（滤过）$A_\infty$ 双模同态之间的同伦。也就
是说，我们将定义（滤过）$A_\infty$ 双模同态之间的同伦、$A_\infty$ 双模之间
的同伦等价等概念。

第 5.3 节给出构造的几何部分。我们将为第 3.7 节中与拉格朗日子流形相对 spin
对 $(L^{(0)},L^{(1)})$ 相伴的滤过 $A_\infty$ 双模构造同伦等价。由此可知，两个
拉格朗日子流形之间的弗洛尔上同调
\[
  HF((L^{(1)},b^{(1)}),(L^{(0)},b^{(0)}))
\]
与构造中涉及的各种选择以及 $L^{(0)},L^{(1)}$ 的哈密顿扰动无关，不过仍依赖
于 $b^{(0)},b^{(1)}$ 的规范等价类。

准确地说，需要以不同方式处理两种不同情形：一种是用同一个辛微分同胚移动
$L^{(0)}$ 和 $L^{(1)}$；另一种是用不同的哈密顿同痕移动二者。在第一种情形中，
取 $\Lambda_{0,\mathrm{nov}}$ 系数的弗洛尔上同调保持不变；在第二种情形中，只有
取 $\Lambda_{\mathrm{nov}}$ 系数的弗洛尔上同调保持不变（见第 1 章的 (Conv.4)）。
第 5.3 节将讨论这一点；第 6 章证明定理 J 时还会再次讨论。更精确地说，有下面
两个定理。

先考虑如下情形。设 $(L^{(0)},L^{(1)})$ 与
$(L^{(0)'},L^{(1)'})$ 分别是辛流形 $M$ 与 $M'$ 中的拉格朗日子流形相对 spin
对。设 $\psi:M\to M'$ 是保持相对 spin 结构的辛微分同胚，且
$\psi(L^{(0)})=L^{(0)'}$、$\psi(L^{(1)})=L^{(1)'}$。固定相容殆复结构族
$J_t,J'_t$ 及其他选择，以定义第 3.7 节中的滤过 $A_\infty$ 双模
\[
  C(L^{(1)},L^{(0)};\Lambda_{0,\mathrm{nov}}),
  \qquad
  C(L^{(1)'},L^{(0)'};\Lambda_{0,\mathrm{nov}}).
\]

\begin{theorem}{4.1.4}\label{thm:4-1-4}
设
\[
  b^{(0)}\in\widehat{\mathcal M}(L^{(0)}),\quad
  b^{(1)}\in\widehat{\mathcal M}(L^{(1)}),\quad
  b^{(0)'}\in\widehat{\mathcal M}(L^{(0)'}),\quad
  b^{(1)'}\in\widehat{\mathcal M}(L^{(1)'}).
\]
假定
$\psi_*([b^{(0)}])=[b^{(0)'}]$ 且
$\psi_*([b^{(1)}])=[b^{(1)'}]$。那么存在同构
{\small
\[
  \psi_*:\quad
  HF((L^{(1)},b^{(1)}),(L^{(0)},b^{(0)});\Lambda_{0,\mathrm{nov}})
  \xrightarrow{\ \cong\ }
  HF((L^{(1)'},b^{(1)'}),(L^{(0)'},b^{(0)'});\Lambda_{0,\mathrm{nov}}).
\]
}
进一步，若 $\{\psi_s\}_s$（$0\le s\le1$）是辛微分同胚的同痕，且每个 $s$
都满足 $\psi_s(L^{(0)})=L^{(0)'}$、$\psi_s(L^{(1)})=L^{(1)'}$，则
$\psi_{0*}=\psi_{1*}$。
\end{theorem}

在定理 4.1.4 中取 $M=M'$ 且 $L^{(i)}=L^{(i)'}$（$i=0,1$），就说明弗洛尔
上同调与定义中涉及的各种选择无关。

例如，考虑 $\dim_{\mathbb R}M=4$、$c_1(M)=0$，且 $L^{(0)},L^{(1)}$ 的马斯洛夫
指标均为零的情形。对 $M$ 上一般的相容殆复结构，不存在以 $L^{(0)}$ 或
$L^{(1)}$ 为边界的伪全纯圆盘，因而我们的滤过 $A_\infty$ 结构等于经典结构；
特别地，可以直接取 $b_0=b_1=0$ 为定界余链。然而，弗洛尔上同调
\[
  HF((L^{(1)},0),(L^{(0)},0);\Lambda_{0,\mathrm{nov}})
\]
依赖于殆复结构的选择。利用定理 4.1.1、4.1.4，可以如下解释这一现象。把殆复
结构由 $J$ 改为 $J'$ 时，得到滤过 $A_\infty$ 同态
\[
  \mathfrak f:
  (C(L^{(i)};\Lambda_{0,\mathrm{nov}}),\mathfrak m)
  \longrightarrow
  (C(L^{(i)};\Lambda_{0,\mathrm{nov}}),\mathfrak m).
\]
注意，因为两端的量子效应都为零，这两个滤过 $A_\infty$ 代数本身并不依赖于
$J$；可是诱导映射
$\mathfrak f_*:\widehat{\mathcal M}(L^{(i)})\to
\widehat{\mathcal M}(L^{(i)})$ 一般并不平凡。因此，由定理 4.1.4 得到同构
\[
  HF((L^{(1)},0),(L^{(0)},0);\Lambda_{0,\mathrm{nov}})
  \cong
  HF((L^{(1)},\mathfrak f_*(0)),
     (L^{(0)},\mathfrak f_*(0));\Lambda_{0,\mathrm{nov}}).
\]
这就是本节开头所说的穿墙现象之一例。

定理 4.1.4 由第 5.3.1 小节的定理 5.3.1 配合第 5.3.7 小节所述的代数论证推出。
下面考虑另一种情形。设 $(L^{(0)},L^{(1)})$ 与
$(L^{(0)'},L^{(1)'})$ 是辛流形 $M$ 中的拉格朗日子流形相对 spin 对；又设
$\{\psi_\rho^{(0)}\}_\rho$、$\{\psi_\rho^{(1)}\}_\rho$（$0\le\rho\le1$）是哈密顿
同痕，满足
\[
  \psi_0^{(0)}=\psi_0^{(1)}=\operatorname{id},\qquad
  \psi_1^{(0)}(L^{(0)})=L^{(0)'},\qquad
  \psi_1^{(1)}(L^{(1)})=L^{(1)'}.
\]

\begin{theorem}{4.1.5}\label{thm:4-1-5}
设
\[
  b^{(0)}\in\widehat{\mathcal M}(L^{(0)}),\quad
  b^{(1)}\in\widehat{\mathcal M}(L^{(1)}),\quad
  b^{(0)'}\in\widehat{\mathcal M}(L^{(0)'}),\quad
  b^{(1)'}\in\widehat{\mathcal M}(L^{(1)'}),
\]
并假定
\[
  \psi_{1*}^{(0)}([b^{(0)}])=[b^{(0)'}],\qquad
  \psi_{1*}^{(1)}([b^{(1)}])=[b^{(1)'}].
\]
这里改取 $\Lambda_{\mathrm{nov}}$ 而非 $\Lambda_{0,\mathrm{nov}}$ 系数的弗洛尔上同调。
那么存在同构
\[
  \begin{aligned}
  &\bigl(\{\psi_\rho^{(1)}\}_\rho,
         \{\psi_\rho^{(0)}\}_\rho\bigr)_*:\\[-0.25ex]
  &\quad HF((L^{(1)},b^{(1)}),(L^{(0)},b^{(0)});\Lambda_{\mathrm{nov}})
  \xrightarrow{\ \cong\ }\\[-0.25ex]
  &\qquad HF((L^{(1)'},b^{(1)'}),(L^{(0)'},b^{(0)'});\Lambda_{\mathrm{nov}}).
  \end{aligned}
\]
进一步，若 $\{\psi_{\rho,s}^{(i)}\}_{\rho,s}$（$i=0,1$）是哈密顿同痕的双参数
族，且对每个 $s$ 都有
$\psi_{0,s}^{(i)}=\operatorname{id}$、
$\psi_{1,s}^{(i)}(L^{(i)})=L^{(i)'}$，则
\[
  \bigl(\{\psi_{\rho,0}^{(1)}\}_\rho,
        \{\psi_{\rho,0}^{(0)}\}_\rho\bigr)_*
  =
  \bigl(\{\psi_{\rho,1}^{(1)}\}_\rho,
        \{\psi_{\rho,1}^{(0)}\}_\rho\bigr)_*.
\]
\end{theorem}
\liftIndexRangeEnd{I140}

这里限于哈密顿同痕，是因为第 5.3.5 小节需要证明能量估计。

定理 4.1.1--4.1.5 还有包含幺元（或同伦幺元）的版本，也有包含第 3.8 节所引
入的 $H^2(M;\Lambda^+_{0,\mathrm{nov}})$ 中元素所给形变的版本。我们将在证明
这些版本的第 5.3.7 小节中陈述它们；定理 G 的 (G.3) 随即由此推出。

定理 A、F 的证明将在第 5.4 节完成。本节主体继续讨论问题的代数方面；其中的
主要结果是定理 W（即定理 5.4.2）和定理 5.4.18，它们都在第 5.4.4 小节证明。
第 5.4 节的另一目的，是把 $\mathcal M(L)$ 描述为某个形式映射（仓西映射）的
零点集；也就是说，我们将在那里证明定理 M。我们还把弗洛尔上同调看成形式概
形 $\mathcal M(L^{(1)})\times\mathcal M(L^{(0)})$ 上凝聚层导出范畴中的对象，
并在第 5.4.6 小节证明定理 N。此外，我们将证明：在对 $L$ 的马斯洛夫类所加的
某些假设下，可以在比 $\Lambda_{0,\mathrm{nov}}$ 略小的环上定义弗洛尔上同调；
特别是，当 $L$ 的马斯洛夫类平凡时，可以在第 1.4 小节定义的 $\Lambda_0$ 上
定义它。第 5.4.6 小节将利用这一点证明推论 O、P。

从现在起，若 $\psi:X\to Y$ 是微分同胚，便约定
\[
  \psi_*=(\psi^{-1})^*,
\]
其中 $(\psi^{-1})^*:H^*(X)\to H^*(Y)$ 是由 $\psi^{-1}$ 诱导的映射。
